\section{Discussion}
\label{sec:discussion}

\paragraph{Did \tool find a new bug?} No. The two potential-recurrence sites
on master (\texttt{X509\_PUBKEY\_set}, \texttt{ossl\_ml\_kem\_i2d\_pubkey};
Observation~\ref{obs:rq1-mastertriage}) share one trigger: the tracked
object's state must change between two reads that a single straight-line
call treats as equal, which only an unsynchronized second thread racing
the same object can cause. Re-checking against the pinned source
(\path{results/reachability_rq1_candidates.md}) also corrected our original
note, which blamed a lock/atomic-op failure in \texttt{EVP\_PKEY\_up\_ref};
every platform variant of \texttt{CRYPTO\_UP\_REF} here returns
\texttt{true}, so that path is dead code. No in-tree caller creates the
precondition, and we built no multithreaded fault-injection test, so both
remain candidates, not findings.

\paragraph{Is this just matching one patch?} The rule names no OpenSSL
function, yet detects the unrelated \texttt{d2217c88df} fix unchanged
(Observation~\ref{obs:rq1-conformance}) and a sibling site fixed later
(Observation~\ref{obs:rq3-timetravel}) -- two data points, not a rate.
Independently of the tool, a commit-message search of OpenSSL's history
(\path{results/historical_recurrence_search.md}) found ten further,
independently-authored fixes of this shape across nine files (2023 to six
weeks before our snapshot), mostly struct-field and local-variable resets
that neither engine tracks. This corroborates the convention's generality
without being a tool finding.

\paragraph{Why a custom engine?} Table~\ref{tab:rq2-comparison} is the
attempt to use existing tools: with stock configurations none separates
any of the three CVEs from its fix (Observation~\ref{obs:rq2-generic-miss}),
because their rule sets do not encode this project's conventions.

\paragraph{False positives and scope.} Across all six snapshots, 1 of 18
flagged Protocol~B sites is a confirmed false positive (\texttt{load\_common},
which releases only a structural reference); two more are correctly
suppressed. For CVE-2023-0215, the Table~\ref{tab:rq2-comparison} distinction comes
from the Part~1 baseline, not \tool; for CVE-2024-4741 no tool, including
the baseline, distinguishes the corrected pair. Protocol~A covers CVE-2023-0215's own pair but its whole-tree
run is 0 for 11 (Section~\ref{sec:rq5}); CVE-2024-4741 needs a state-guard
rule neither shape expresses. Finally, the 91.7\% master conformance rate
is not a precision number: it holds only under the name-based release
heuristic, and sites that heuristic cannot see are absent from both
numerator and denominator.
